\subsection{Screening and reuse}
\label{app:screened-release}
This diagnostic was declared after the negative real-data study. It
does not replace that frozen evaluation. At minimum gain $1/200$, all
nine public problems fail the information-value screen: eight have
zero nonprivate value beyond reuse; TLC at $n=4$ permits only
$8.6893\times10^{-6}$. Enumerating all 256 old-output maps and an
attaining state-dependent oracle independently confirms these values.
The post-LP screen gives the same decisions.

One warm call followed by five timed calls gives median screen times
2.52--28.24\,ms on already constructed public problems, excluding public
fitting, domain construction and JSON loading. This avoids an LP, not
finite-state enumeration. The bound concerns the declared public prior
and output family, not unknown population accuracy.

Three separate $n=4$ workflows use actual calibration records, the four
binary-feature heads and a reuse certificate constructed directly from
the public problem, with LP calls disabled. Registration and exact
validation precede the old draw. A prospective charge and cap of
$78/125=.624$ cover old-law odds $97/52$; an odds-$3$ plan is rejected.
Committed model bytes survive cold-cache ledger reopening, retries and
delivery to another recipient. Another charged history is refused at
the exhausted cap. Revocation blocks subsequent delivery but retains
the charge. An identical-channel control retains its earlier $1.1$
reservation after revocation: tighter prospective accounting neither
improves an unchanged channel's statistical privacy nor refunds history.
These local checks establish no private-learning gain or agency deployment.
